\documentclass[12pt,a4paper]{article}
\usepackage[utf8]{vntex, inputenc}
\usepackage{amsmath, amssymb, anyfontsize, array, color,  enumerate, enumitem, fancyhdr, float, graphicx, hyperref, lastpage, listings, mathtools, mips, natbib, pifont, tikz}
\usepackage[open,openlevel=0]{bookmark}
\usepackage[framemethod=tikz]{mdframed}

\usetikzlibrary{calc}
\newcommand\HRule{\rule{12cm}{1pt}}

\pagestyle{fancy}
\fancyhead{} % clear all header fields
\fancyhead[L]{
	\color{blue}
	\begin{tabular}{rl}
		\begin{picture}(25,15)(0,0)
		\put(0,-8){\includegraphics[width=8mm, height=8mm]{BK.png}}
		\end{picture}
		\begin{tabular}{l}
			\textbf{\bf \ttfamily Ho Chi Minh University of Technology}\\
			\textbf{\bf \ttfamily Faculty of Computer Science \& Engineering}
		\end{tabular}
	\end{tabular}
}
\fancyhead[R]{
	\begin{tabular}{l}
		\tiny \bf \\
		\tiny \bf
	\end{tabular}
}
\fancyfoot{} % clear all footer fields
\fancyfoot[L]{\scriptsize \ttfamily Computer Architecture Week 3}
\fancyfoot[R]{\scriptsize \ttfamily Page {\thepage}/\pageref{LastPage}}
\renewcommand{\headrulewidth}{0.3pt}
\renewcommand{\footrulewidth}{0.3pt}

\definecolor{dkgreen}{rgb}{0,0.6,0}
\definecolor{gray}{rgb}{0.5,0.5,0.5}
\definecolor{mauve}{rgb}{0.58,0,0.82}

\lstset{frame=tb,
	language=[MIPS]Assembler,
	aboveskip=3mm,
	belowskip=3mm,
	showstringspaces=false,
	columns=flexible,
	basicstyle={\small\ttfamily},
	numbers=none,
	numberstyle=\tiny\color{gray},
	keywordstyle=\color{blue},
	commentstyle=\color{dkgreen},
	stringstyle=\color{mauve},
	breaklines=true,
	breakatwhitespace=true,
	tabsize=3,
	numbers=left,
	stepnumber=1,
	numbersep=1pt,
	firstnumber=1,
	numberfirstline=true
}

\begin{document}
\begin{titlepage}
  \begin{tikzpicture}[remember picture, overlay]
    \draw[line width = 2pt] ($(current page.north west) + (1in,-1in)$) rectangle ($(current page.south east) + (-1in,1in)$);
  \end{tikzpicture}

  \begin{center}

    % Upper part of the page
    \textbf{\fontsize{12pt}{1pt}\selectfont HO CHI MINH CITY UNIVERSITY OF TECHNOLOGY}\\[0.5cm]
    {\fontsize{13pt}{1pt}\selectfont Faculty of Computer Science \& Engineering}\\[0.5cm]
    \begin{figure}[H]
      \centering
      \includegraphics[width=1.7in,height=1.7in]{BK.png}
    \end{figure}

    % Title
    \HRule\\[0.5cm]
    { \textbf{\fontsize{25pt}{1pt}\selectfont Computer Architecture Week 3}}\\[0.4cm]

    \HRule\\[0.8cm]
    \begin{minipage}{0.545\textwidth}
      \begin{flushleft}
        \textbf{Authors:}\\
        \begin{tabular}{l l}
          Lưu Nguyễn Hoàng Minh & 1952845 \\
          Vũ Anh Nhi            & 1952380 \\
          Nguyễn Chính Khôi     & 1952793 \\
          Nguyễn Hoàng          & 1952255 \\
        \end{tabular}
      \end{flushleft}
    \end{minipage}
    \begin{minipage}{0.4\textwidth}
      \begin{flushright}
        \textbf{Class:}\\
        CO2008\_CC02\\
        \textbf{Lecturer:}\\
        Kiều Đỗ Nguyên Bình\\

      \end{flushright}
    \end{minipage}

    \vfill

    % Bottom of the page
    \vspace{2cm}
    {\large} %{\large \today}
  \end{center}
\end{titlepage}


\pdfbookmark[0]{Question 1}{ques1}
\section*{Question 1:}
Write a MIPS program with the following requirements:
\begin{enumerate}[label=\alph*.]
  \item Declare an integer array with 10 synthetic data elements.
  \item Calculate the sum of all elements.
  \item Print the result to the terminal.
\end{enumerate}

\begin{mdframed}[hidealllines=true,backgroundcolor=magenta!10]
  \begin{lstlisting}

      .data:
      X:      .word 0, 1, 2, 3, 4, 5, 6, 7, 8, 9
      msg1:   .asciiz "The sum of all elements is "

      # Start program code
      .text:

      .globl main
  main:
      la      $s7, X              # Address of X
      add     $s0, $zero, $zero   # Counter
      add     $s1, $zero, $zero   # Sum
      addi    $t7, $zero, 40      # Limiter
  Sum:
      add     $t0, $s0, $s7       # Address of element
      lw      $t1, 0($t0)
      add     $s1, $s1, $t1       # Summer
      addi    $s0, $s0, 4         # i = i + 1
      beq     $s0, $t7, SExit
      j       Sum

  SExit:
      # Print string msg1
      li      $v0, 4
      la      $a0, msg1
      syscall

      # Print sum
      li      $v0, 1
      move    $a0, $s1
      syscall

      # Exit
      li      $v0, 10
      syscall

  \end{lstlisting}
\end{mdframed}

\pdfbookmark[0]{Question 2}{ques2}
\section*{Question 2:}
Write a MIPS program with the following requirements:
\begin{enumerate}[label=\alph*.]
  \item Declare an integer array with 10 synthetic data elements.
  \item Calculate the sum of all odd elements (a[1], a[3],\dots).
  \item Calculate the sum of all even elements (a[0], a[2],\dots).
  \item Print the results to the terminal.
\end{enumerate}

\begin{mdframed}[hidealllines=true,backgroundcolor=magenta!10]
  \begin{lstlisting}

      .data:
      X:      .word 0, 1, 2, 3, 4, 5, 6, 7, 8, 9
      msg1:   .asciiz "The sum of all elements of even indexes is "
      msg2:   .asciiz "The sum of all elements of odd indexes is "
      lf:     .asciiz "\n"

      # Start program code
      .text:

      .globl main
  main:
      la      $s7, X              # Address of X
      add     $s0, $zero, $zero   # Counter
      add     $s1, $zero, $zero   # Sum of even
      add     $s2, $zero, $zero   # Sum of odd
      addi    $t7, $zero, 40      # Limiter

  Sum:
      add     $t0, $s0, $s7       # Address of element
      lw      $t1, 0($t0)
      add     $s1, $s1, $t1       # Even summer
      addi    $s0, $s0, 4         # i = i + 1
      beq     $s0, $t7, SExit

      add     $t0, $s0, $s7       # Address of element
      lw      $t1, 0($t0)
      add     $s2, $s2, $t1       # Odd summer
      addi    $s0, $s0, 4         # i = i + 1
      beq     $s0, $t7, SExit

      j       Sum

  SExit:
      # Print string msg1
      li      $v0, 4
      la      $a0, msg1
      syscall

      # Print sum of odd
      li      $v0, 1
      move    $a0, $s1
      syscall

      # Print string new line
      li      $v0, 4
      la      $a0, lf
      syscall

      # Print string msg2
      li      $v0, 4
      la      $a0, msg2
      syscall

      # Print sum of even
      li      $v0, 1
      move    $a0, $s2
      syscall

      # Exit
      li      $v0, 10
      syscall

  \end{lstlisting}
\end{mdframed}

\pdfbookmark[0]{Question 3}{ques3}
\section*{Question 3:}
Write a MIPS proram that receives 10 integer numbers from users through the terminal and store these numbers into an array. Print the sum of all array elements to the terminal.

\begin{mdframed}[hidealllines=true,backgroundcolor=magenta!10]
  \begin{lstlisting}

      .data:
      X: .space 40 # Reserve 40 bytes
      msg1: .asciiz "Enter an integer: "
      msg2: .asciiz "The sum of all elements is "

      # Start program code
      .text:

      .globl main
  main:
      la      $s7, X              # Address of X
      add     $s0, $zero, $zero   # Counter
      add     $s1, $zero, $zero   # Sum
      addi    $t7, $zero, 40      # Limiter
  Input:
      add     $t0, $s0, $s7       # Address of element

      # Print string msg1
      li      $v0, 4
      la      $a0, msg1
      syscall

      # Get input and save
      li      $v0, 5
      syscall
      move    $t1, $v0

      sw      $t1, 0($t0)
      addi    $s0, $s0, 4         # i = i + 1
      beq     $s0, $t7, IExit
      j       Input

  IExit:
      add     $s0, $zero, $zero   # Counter
      add     $s1, $zero, $zero   # Sum
  Sum:
      add     $t0, $s0, $s7       # Address of element
      lw      $t1, 0($t0)
      add     $s1, $s1, $t1       # Summer
      addi    $s0, $s0, 4         # i = i + 1
      beq     $s0, $t7, SExit
      j       Sum

  SExit:
      # Print string msg2
      li      $v0, 4
      la      $a0, msg2
      syscall

      # Print sum
      li      $v0, 1
      move    $a0, $s1
      syscall

      # Exit
      li      $v0, 10
      syscall

    \end{lstlisting}
\end{mdframed}


\pdfbookmark[0]{Question 4}{ques4}
\section*{Question 4:}
Write a MIPS program with the following requirements:
\begin{enumerate}[label=\alph*.]
  \item Declare an integer array with 10 synthetic data elements.
  \item Print a sentence to terminal to request an integer number that is greater than 0 and less than 10 (assume that user strictly follow this rule).
  \item Print value of the element at index collected from the previous step.
\end{enumerate}

\begin{mdframed}[hidealllines=true,backgroundcolor=magenta!10]
  \begin{lstlisting}

      .data:
      X: .word 0, 11, 22, 33, 44, 55, 66, 77, 88, 99
      msg1: .asciiz "Enter an integer between 0 and 10: "
      msg2: .asciiz "Data at index is "

      # Start program code
      .text:

      .globl main
  main:
      la      $s7, X              # Address of X
      add     $s0, $zero, $zero   # Counter
      addi    $t6, $zero, 4       # Sum limiter
      addi    $t7, $zero, 40      # Limiter

      # Print string msg1
      li      $v0, 4
      la      $a0, msg1
      syscall

      # Get input and save
      li      $v0, 5
      syscall
      move    $t1, $v0
      mul     $t6, $t1, $t6

  Search:
      add     $t0, $s0, $s7       # Address of element
      lw      $t1, 0($t0)
      beq     $s0, $t6, SExit
      beq     $s0, $t7, SExit
      addi    $s0, $s0, 4         # i = i + 1
      j       Search

  SExit:
      # Print string msg2
      li      $v0, 4
      la      $a0, msg2
      syscall

      # Print sum
      li      $v0, 1
      move    $a0, $t1
      syscall

      # Exit
      li      $v0, 10
      syscall

  \end{lstlisting}
\end{mdframed}


\pdfbookmark[0]{Question 5}{ques5}
\section*{Question 5:}

Write a MIPS program that reverses an 10 elements integer array. For example, the array
initially stores 1, 3, 5, 7, 9, 11, 13, 15, 17, 19, the program will change the array to be 19, 17, 15, 13, 11, 9,
7, 5, 3, 1.

\begin{mdframed}[hidealllines=true,backgroundcolor=magenta!10]
  \begin{lstlisting}

    .data
    arr: .word 34, 23, 15, 82, 38, 56, 48, 93, 49, 11
    newarr: .word 0:10
    space: .asciiz " "
    endl: .asciiz "\n"

    .text
    .globl main

  main:
    la $s0, arr
    add $s1, $zero, $zero

  L1: #print the original array
    lw $t1, 0($s0)
    li	$v0,1
    move	$a0, $t1
    syscall

    addi $s0, $s0, 4
    addi $s1, $s1, 1

    li	$v0, 4
    la	$a0, space
    syscall

    beq $s1, 10, r1
    j L1

  r1:
    li	$v0, 4
    la	$a0, endl
    la $s3, newarr
    syscall
    addi $s0, $s0, -4

  rev: # reverse the array
    lw $t2, 0($s0)
    sw $t2, 0($s3)
    addi $s0, $s0, -4
    addi $s3, $s3, 4
    addi $s1, $s1, -1

    beq $s1, 0, print
    j rev

  print:
    la $s3, newarr
    addi $s1, $zero, 0
    j print2

  print2: # print the reverse array
    lw $t1, 0($s3)
    li	$v0,1
    move	$a0, $t1
    syscall

    addi $s3, $s3, 4
    addi $s1, $s1, 1

    li	$v0, 4
    la	$a0, space
    syscall

    beq $s1, 10, Exit
    j print2

  Exit:
    li $v0, 10
    syscall

    \end{lstlisting}
\end{mdframed}

\end{document}